\documentclass[12pt]{article}
\usepackage[utf8]{inputenc}
\usepackage[english]{babel}
\usepackage{amssymb}
\usepackage{stmaryrd}
\usepackage{enumitem}
\usepackage{amsmath}
\usepackage{array}
\usepackage[left=2cm,right=2cm,top=2cm,bottom=2cm]{geometry}
\usepackage[T1]{fontenc}
\setlength\parindent{0pt}
\usepackage{graphicx}
\usepackage[ddmmyyyy]{datetime}

\usepackage[backend=biber]{biblatex}
\addbibresource{ref.bib}

\usepackage{booktabs}
\usepackage{float}
\usepackage{forest}
\usepackage{setspace}
\usepackage{hyperref}
\hypersetup{
    colorlinks=true,
    linkcolor=blue,
    filecolor=magenta,      
    urlcolor=cyan,
}
\usepackage{easytable}
\usepackage{adjustbox}
\usepackage{listings}
\lstdefinelanguage{Matlab}%
  {morekeywords={matlab2tikz},
   morekeywords=[2]{1}, 
   morekeywords=[3]{xlim,ylim,clim,iter},
   morekeywords=[4]{x,y,scatter,contour,imagesc,addplot3},
   morecomment=[l][\color{blue}]{...},%
   morestring=[m]',%
   morestring=[m]"%
  }[keywords,comments,strings]

\lstset{%
  language=Matlab,%
  breaklines=true,%
  morekeywords={matlab2tikz},
  keywordstyle=\color{blue},%
  morekeywords=[2]{1}, keywordstyle=[2]{\color{black}},
  identifierstyle=\color{black},%
  stringstyle=\color{mylilas},
  commentstyle=\color{mygreen},%
  showstringspaces=false,
  numbers=left,%
  numberstyle={\tiny \color{black}},
  numbersep=9pt,
  emph=[1]{for,end,break},emphstyle=[1]\color{red},
}

\hypersetup{
    colorlinks=true,
    linkcolor=black,  % Changer la couleur ici
    filecolor=magenta,      
    urlcolor=cyan,
}

\usepackage{tikz}
\usetikzlibrary{calc}

\renewcommand{\dateseparator}{/}

\usepackage{fancyhdr}
\usepackage{lastpage}

\pagestyle{fancy}
\renewcommand\headrulewidth{1pt}
\renewcommand\footrulewidth{1pt}
\geometry{headsep=1.1cm}

\newtheorem{theorem}{Théorème}[subsection]
\newtheorem{corollary}{Corollaire}[theorem]
\newtheorem{lemma}[theorem]{Lemme}
\newtheorem{definition}{Définition}[subsection]
\newtheorem{proposition}{Proposition}[section]

\fancyhead[L]{\includegraphics[width=0.1\columnwidth]{./logo}~}
\fancyfoot[L]{\textsc{TD3 Optimisation}}
\fancyhead[R]{\textsc{Adam YOUNSI, Nabil ALAMI, Edward LUCYSZYN}}
\fancyfoot[C]{\thepage/\pageref{LastPage}}
\fancyfoot[R]{\textsc{\today}}

\usepackage{xcolor}
\usepackage{piton}
\usepackage{tcolorbox}
\newcommand{\Tr}{\operatorname{Tr}}
\newcommand{\Arg}{\operatorname{Arg}}

\usepackage{amsmath}
\DeclareMathOperator*{\argmax}{argmax}
\DeclareMathOperator*{\argmin}{argmin}

%====================== INFORMATION ET REGLES ======================

%rajouter les numérotation pour les \paragraphe et \subparagraphe

\setcounter{tocdepth}{4}
\setcounter{secnumdepth}{0}


%======================== DEBUT DU DOCUMENT ========================

\begin{document}

% \NewPitonEnvironment{Python}{O{}}{\PitonOptions{#1}}{}
\PitonOptions{break-lines, indent-broken-lines}

\NewPitonEnvironment{Python}{}
  {\begin{tcolorbox}[breakable]}
  {\end{tcolorbox}}


%régler l'espacement entre les lignes
\newcommand{\HRule}{\rule{\linewidth}{0.5mm}}

%page de garde
\begin{titlepage}
\begin{center}

% Upper part of the page. The '~' is needed because only works if a paragraph has started.

\LARGE \textsc{2EL1740: Algèbre \& Cryptologie}

\vspace{0.2cm}

\Large \textsc{CentraleSupélec - 2A}

\vspace{0.3cm}

% Title
\HRule \\[0.4cm]

{\huge \bfseries Challenge 4\\
[0.2cm]}

\HRule \\[0.4cm]

\vspace{2cm}

\textsc{\today}

\vspace{2cm}

\includegraphics[width=0.4\columnwidth]{./logo}~\\[3cm]

% Author and supervisor
\begin{minipage}{0.4\textwidth}
\begin{spacing}{1.125}
\begin{center}
    Raphaël \testsc{PAIN DIT HERMIER}\\
    Alexis \testsc{LOMBARD-GAILLARD}\\
    Edward \testsc{LUCYSZYN}
\end{center}
\end{spacing}
\end{minipage}

\vfill

\end{center}
\end{titlepage}

%page blanche
%\newpage
%~
%ne pas numéroter cette page

\tableofcontents
\setcounter{page}{1}
%ne pas numéroter le sommaire

%\newpage

%espacement entre les lignes d'un tableau
\renewcommand{\arraystretch}{1.5}

%====================== INCLUSION DES PARTIES ======================


%recommencer la numérotation des pages à "1"

\newpage
\section{One-Dimensional Schrödinger Equation}

\noindent The goal is to establish the necessary mathematical foundations to better understand and assimilate the concepts of wave function and the resolution of the Schrödinger equation.

\subsection{First Order Equations}
\noindent Consider the following time-dependent equation:
\begin{equation}
    \forall t \in \mathbb{R}_{+}, \quad \frac{df}{dt}(t) = af(t)
\end{equation}
where $\displaystyle a=-\frac{1}{\tau}$ is a real constant ($\tau >0$).\\ \\

\noindent \textbf{1)} Find the set of solutions to the differential equation. What is the dimension of $\tau$?\\

\begin{breakbox}
\noindent This is a first-order differential equation. We know that $f$ has the form: $$\boxed{\forall t \in \mathbb{R}_{+}, \quad f(t)=Ae^{at}}$$ 
where $A$ is a constant (in $\mathbb{C}$ generally) that depends on an initial condition (or normalization condition). $\tau$ has the dimension of time. It is often referred to as the \textbf{characteristic time}.
\end{breakbox}

\subsection{Second Order Equations}

\begin{equation}
    \frac{d^2\phi}{dx^2}(x) = b\phi(x)
    \label{equ2}
\end{equation}

\noindent
where $b$ is a nonzero real constant.\\ \\

\noindent \textbf{2)} Find the set of solutions to the differential equation. Solutions can be sought in exponential form, distinguishing different cases based on the sign of $b$.\\

\begin{breakbox}

\noindent There are several methods to solve this second-order differential equation.\\ \\
\noindent \textbf{Method 1: Seeking solutions of the form $\phi (x) = Ae^{cx}, c \in \mathbb{R}$.}\\
We seek solutions in exponential form (this point is admitted). So, $$\phi (x) = Ae^{cx}$$ with $c$ to be determined.
We substitute this solution into the differential equation (2): $$c^2Ae^{cx} = b Ae^{cx}.$$
By simplifying with $\phi (x)= Ae^{cx}$ ($A \neq 0$), we obtain: $c^2 = b$. \\ \\Here we distinguish two cases:

\begin{itemize}
    \item If $b<0$: we set $b=(ik)^2$, with $k > 0$. Then we have the equality $c^2 = (ik)^2$, 
resulting in $c = \pm ik$. There are thus two forms of exponential solutions, one in $ikx$ and the other in $-ikx$. Since these two solutions are independent, the general solution for $\phi$ is: $$\boxed{\phi(x) = Ae^{ikx} + Be^{-ikx}.}$$
    \item If $b>0$: we set $b=k^2$. Then we find $c = \pm k$.
As before, we find two forms of independent exponential solutions, and the general solution is: $$\boxed{\phi = Ae^{kx} + Be^{-kx}.}$$
\end{itemize}
\textbf{Method 2: Passing through the characteristic equation.}\\
We write the characteristic equation (or polynomial) in $X$ associated with the differential equation (2), which is:
\begin{equation}
    X^2 - b = 0.
    \label{eq3}
\end{equation}
We also know that the solutions to (\ref{equ2}) are of the form $\psi(x) = Ae^{r_1x}+Be^{r_2x}$ where $r_1, r_2 \in \mathbb{C}$ are the complex roots of equation (\ref{eq3}).\\
We then distinguish the cases $b >0$ and $b<0$ as detailed in Method 1.

\end{breakbox}

\subsection{Variable Separation and Schrödinger}

\noindent Let $\psi$ be a function of two variables that satisfies the one-dimensional Schrödinger equation: 
\begin{equation}
    \forall (x, t) \in \mathbb{R} \times \mathbb{R}_{+}, - \frac{\hbar ^2}{2m}\frac{\partial^2\psi(x,t)}{\partial x^2} = i\hbar\frac{\partial\psi(x,t)}{\partial t}.
\end{equation}
\noindent Furthermore, suppose that $\psi$ can be written in the form: 
\begin{equation}
    \psi(x,t) = \phi(x)f(t).
\end{equation}


\noindent \textbf{3.a)} Simplify the equation to obtain a term dependent only on $x$ and another term dependent only on $t$. What can be concluded from this equality? \\

\begin{breakbox}
\noindent We first notice that $$\frac{\partial^2\psi(x,t)}{\partial x^2} = \frac{\partial^2\phi(x)f(t)}{\partial x^2} = f(t)\phi^{''}(x).$$
Likewise:  $$\frac{\partial\psi(x,t)}{\partial t} = \phi(x)\dot{f}(t).$$
Knowing this, we rewrite the Schrödinger equation: $$- \frac{\hbar ^2}{2m}f(t)\phi^{''}(x) = i\hbar\phi(x)\dot{f}(t).$$
Thus finally:
$$\boxed{-\frac{\hbar ^2}{2m}\frac{\phi^{''}(x)}{\phi(x)} = i\hbar\frac{\dot{f}(t)}{f(t)}.}$$\\

\noindent We obtain an equality between an expression dependent only on $x$ and another dependent only on $t$.
Therefore, we conclude that the equality must necessarily be equal to a constant, which we denote $E \in \mathbb{R}$.
\end{breakbox} 

\medskip

\noindent \textbf{3.b)} Deduce the differential equations satisfied by $\phi$ and $f$. Provide the dimension and physical meaning of the introduced constant.\\

\begin{breakbox}

\noindent The constant $E$ has dimensions of energy, which can be seen by examining the dimension of one of the two terms. For example: 
$$[E]=[\frac{\hbar ^2}{2m}\frac{\phi^{''}(x)}{\phi(x)}\big] = [\frac{\hbar ^2}{2m}] L^{-2} = [\frac{\hbar^2 k^2}{2m}] = [\frac{p^2}{2m}] = Energy$$ \\ \\
We then obtain the following two differential equations:

\begin{equation}
    \boxed{\dot{f}(t) = \frac{E}{i\hbar}f(t);}
    \label{a1}
\end{equation}

\begin{equation}
    \boxed{\phi^{''}(x) = -\frac{2mE}{\hbar^2}\phi(x).}
    \label{a2}
\end{equation}
\end{breakbox}

\medskip

\noindent \textbf{3.c)} Then give the general form of the solutions for $\psi(x,t)$ depending on the sign of the introduced constant.\\

\begin{breakbox}
\noindent Here it is a matter of recognizing the forms of the differential equations from questions 1.1 for equation (\ref{a1}), and 1.2 for equation (\ref{a2}). Thus $$f(t)= Ae^{-i \frac{Et}{\hbar}}.$$\\ \\

\begin{itemize}
    \item If $E$ is negative, (\ref{a2}) becomes $$\phi^{''}(x) = \frac{2m|E|}{\hbar^2}\phi(x)$$ and according to 1.2:
    $$\phi(x) = Ae^{kx} + Be^{-kx}\text{, where } k^2 = -\frac{2mE}{\hbar^2}.$$ Finally:
    $$\boxed{\psi(x,t) = (Ae^{kx} + Be^{-kx})e^{\frac{-iEt}{\hbar}}.}$$
    \\
    \item Conversely, if $E$ is positive, then as shown in question 2, the solution will be of the form $\phi(x) = Ae^{ikx} + Be^{-ikx}$, with $\displaystyle k^2 = \frac{2m
    E}{\hbar^2}.$ \\ \\        Finally: $$\boxed{\psi(x,t) = (Ae^{ikx} + Be^{-ikx})e^{\frac{-iEt}{\hbar}}.}$$
\end{itemize}
\end{breakbox}


% \newpage

% \printbibliography[title={Bibliography}]

\end{document}